% ══════════════════════════════════════════════════════════════════════════════
% PALETA DEL LABORATORIO DE R PARA CIENCIA POLÍTICA — ITAM
% Verde institucional #006847 sobre blanco
%
% ESTE ES EL ÚNICO LUGAR NORMATIVO DONDE VIVEN LOS HEXES.
% Si cambia un color, se cambia aquí y se propaga a mano a los otros dos
% archivos que lo heredan:
%   · estilo/tema_lab.R     → theme_lab() y las escalas de ggplot2
%   · estilo/estilo.scss    → el sitio Quarto
% Los tres tienen que decir lo mismo. Si no, las láminas, los gráficos y el
% sitio dejan de combinar y se nota.
% ══════════════════════════════════════════════════════════════════════════════

\definecolor{verdeITAM}{HTML}{006847}     % acento principal — verde institucional
\definecolor{verdeProfundo}{HTML}{00402C} % títulos y texto enfático
\definecolor{verdeClaro}{HTML}{A8D5BE}    % fondos suaves, resaltados
\definecolor{grisTexto}{HTML}{3D3D3D}     % cuerpo de texto
\definecolor{guinda}{HTML}{7B1113}        % acento secundario, uso escaso

% ── Derivados ────────────────────────────────────────────────────────────────
% No son colores nuevos: son diluciones de los anteriores. Se definen aquí para
% que nadie escriba "verdeClaro!15" a mano en una lámina y el curso termine con
% quince tonos distintos repartidos por ahí.

\colorlet{verdeVelo}{verdeClaro!22!white}   % fondo de bloque y de código
\colorlet{verdeTenue}{verdeClaro!55!white}  % resaltado dentro de un bloque
\colorlet{verdeSombra}{white!30!verdeITAM}  % texto atenuado SOBRE fondo verde
\colorlet{grisVelo}{grisTexto!7!white}      % fondos muy tenues, reglas
\colorlet{grisTenue}{grisTexto!45!white}    % texto terciario, pies
